\question What is the order of growth in time for the following functions? Give
tight bounds in $\Theta$ notation. Assume \lstinline{n} is a non-negative
integer and arithmetic takes constant time. If different inputs follow
different execution paths, describe their runtimes separately.
\begin{parts}

\part
\begin{lstlisting}
def strange_add(n):
    if n == 0:
        return 1
    else:
        return strange_add(n - 1) + strange_add(n - 1)
\end{lstlisting}
\begin{solution}[0.25in]
$\Theta(2^n)$. Every non-base call makes two recursive calls. At depth $k$
there are $2^k$ calls, for depths 0 through $n$. The total is
$1 + 2 + \cdots + 2^n = 2^{n+1} - 1$ calls, with constant work per call.
\end{solution}

\part
\begin{lstlisting}
def stranger_add(n):
    if n < 3:
        return n
    elif n % 3 ==  0:
        return stranger_add(n - 1) + stranger_add(n - 2) + stranger_add(n - 3)
    else:
        return n
\end{lstlisting}
\begin{solution}[0.25in]
$\Theta(n)$ for positive multiples of 3; $\Theta(1)$ otherwise, including
\lstinline{n = 0}.

Only calls whose argument is a positive multiple of 3 recurse. For example,
\lstinline{stranger_add(9)} calls the function on 8, 7, and 6. The calls on
8 and 7 return immediately through the \lstinline{else} branch. The call on
6 calls the function on 5, 4, and 3; again, only the last call continues.
The continuing chain is $9 \to 6 \to 3 \to 0$.

Thus, each recursive stage does constant work and makes two constant-time
calls plus one call on $n-3$: $T(n)=T(n-3)+\Theta(1)$.
There are $n/3$ such stages, giving $\Theta(n)$ time. There are three calls
in the code, but they do not form three continuing recursive branches.
For inputs not divisible by 3, the initial call returns immediately.
\end{solution}
\end{parts}
